% Point to the main file
\documentclass[../../Main.tex]{subfiles}

% ========== Start the document ==========
\begin{document}

\section{Review for Exam 3}

\subsection{Wilson's Theorem}

\begin{theorem}[Wilson's Theorem]
    If $p$ is a prime integer, then

    $$(p - 1)! \equiv p - 1 \pmod{p}$$
    
    or

    $$(p - 1)! \equiv -1 \pmod{p}$$
\end{theorem}

But for composite numbers, we have

\begin{theorem}[Some Unamed Theorem]
    If $n$ is a composite integer greater than 4,

    then

    $$(n-1)! \equiv 0 \pmod{n}$$

    If $n = 4$, then $$3! \equiv 2 \pmod{4}$$
\end{theorem}

\subsection{Fermat's Little Theorem}

\begin{theorem}[Fermat's Little Theorem]
    If $a \in \setZ$ with, 
    
    \begin{enumerate}
        \item $p$ is prime
        \item $\text{gcd}(a, p) = 1$
    \end{enumerate}
 
    $$a^{p - 1} \equiv 1 \pmod{p}$$
\end{theorem}

\subsection{Extended Fermat's Little Theorem}

\begin{theorem}[Extended Fermat's Little Theorem]
    If $a, p, q \in \setZ$ with
    
    \begin{enumerate}
        \item $p$ and $q$ are distinct primes
        \item gcd$(a, pq) = 1$
    \end{enumerate}

    Then

    $$a^{(p-1)(q-1)} \equiv 1 \pmod{pq}$$
\end{theorem}

\subsection{Bezout's Theorem}

\begin{theorem}[B\'{e}zout's Theorem]

    Basically

    $$a, b \in \setZ^+ \implies \exists x, y \in \setZ, ax + by = \text{gcd}(a, b)$$
\end{theorem}

When given a problem that's like find $x, y$ in $ax + by = \text{gcd}(a,b)$ (where they give you $a$ and $b$), just
do Euclidean and stop when you get to the actual GCD then back solve

\subsection{Is a number prime or composite?}

Lowkey just brute force up until 

$$\lceil \sqrt{n} \rceil$$

\subsection{Prime Factorization}

Let's look at how to find the number of divisors in a number $n$

\begin{definition}[Prime Factorization]
    \textbf{Prime factorization} is breaking up a number into prime factors. An example is

    \begin{align*}
        24 &= 6 \cdot 4 \\
        &= (3 \cdot 2)(2 \cdot 2) \\
        &= 3^1 \cdot 2^3
    \end{align*}

    The prime factors are 3 and 2.
\end{definition}

In order to find the number of divisors, we must first break up $n$ into it's prime factors

$$n = p_1^{e_1} \cdot p_2^{e_2} \cdot \ldots \cdot p_k^{e_k}$$

\begin{center}
    $p_k$: Prime factor \\
    $e_k$: Occurences
\end{center}

The number of divisors is then

$$(e_1 + 1)(e_2 + 1) \ldots (e_k + 1)$$

\subsection{Chinese Remainder Theorem}

Given the system: (I'm just gonna use up to 3 for this)

\begin{align*}
    x &\equiv a_1 \pmod{m_1} \\
    x &\equiv a_2 \pmod{m_2} \\
    x &\equiv a_3 \pmod{m_3} \\
\end{align*}

The system has a \textbf{unique solution} if $m_1, m_2, m_3$ are pairwise coprime. If not, then it \textit{could possibly} have a solution but not necessarily

\begin{enumerate}
    \item 
    {
        Expand the first equation

        \begin{align*}
            x &\equiv a_1 \pmod{m_1}
        \end{align*}

        implies

        \begin{align*}
            \Aboxed{x = m_1 k_1 + a_1}, \; k_1 \in \setZ
        \end{align*}
    }

    \item
    {
        Plug that into the second equation

        \begin{align*}
            x &\equiv a_2 \pmod{m_2} \\
            m_1 k_1 + a_1 &\equiv a_2 \pmod{m_2} \\
        \end{align*}

        And solve for $k_1$

        \begin{align*}
            m_1 k_1 + a_1 &\equiv a_2 \pmod{m_2} \\
            m_1 k_1 &\equiv a_2 - a_1 \pmod{m_2} \\
            m_1^{-1} k_1 &\equiv m_1^{-1}(a_2 - a_1) \pmod{m_2} \\
            k_1 &\equiv m_1^{-1}(a_2 - a_1) \pmod{m_2} \\
            \Aboxed{k_1 &\equiv c_1} \pmod{m_2} \tag{$c_1$ is shorthand} \\
        \end{align*}

        \begin{note}
            This isn't a formula, just steps with arbitrary variables. You'll need to find the inverse of $m_1$
        \end{note}

        And expand that last equation

        \begin{align*}
            k_1 &\equiv c_1 \pmod{m_2} \\
        \end{align*}

        implies

        \begin{align*}
            \Aboxed{k_1 &= m_2 k_2 + c_1}, \; k_2 \in \setZ \\
        \end{align*}
    } 

    \item 
    {
        Plug that into our equation for $x$

        \begin{align*}
            x &= m_1 k_1 + a_1 \\
            \Aboxed{x &= m_1 (m_2 k_2 + c_1) + a_1} \\
        \end{align*}

        Plug that into the last equation

        \begin{align*}
            m_1 (m_2 k_2 + c_1) + a_1 &\equiv a_3 \pmod{m_3}
        \end{align*}

        and solve for $k_2 \pmod{m_3}$

        \begin{align*}
            &\vdots \\
            \Aboxed{k_2 &\equiv c_2 \pmod{m_3}} \\
        \end{align*}

        Expand that equation

        \begin{align*}
            \Aboxed{k_2 = m_3 k_3 + c_2}, \; k_3 \in \setZ
        \end{align*}
    }

    \item 
    {
        Finally, plug that last equation into our SECOND equation for $x$

        \begin{align*}
            x &= m_1 (m_2 k_2 + c_1) + a_1 \\
            x &= m_1(m_2(m_3 k_3 + c_2) + c_1) + a_1 \\
            \Aboxed{x &= m_1m_2m_3k_3 + c_3} \\
        \end{align*}

        Now we UN-expand this equation and see that

        \begin{align*}
            \Aboxed{x &\equiv c_3} \pmod{m_1m_2m_3}
        \end{align*}
    }

\end{enumerate}

Basically, starting with the first equation

\begin{enumerate}
    \item \yellowbox{Expand first, solving for $x$}
    \item \yellowbox{Plug $x$ into second}
    \item \greenbox{Solve for $k_1$}
    \item \greenbox{Expand $k_1$}
    \item \greenbox{Plug $k_1$ back into first expanded of $x$ (one with $k_1$)}
    \item Solve for $k_2$
    \item Expand $k_2$
    \item Plug $k_2$ back into second expanded of $x$ (one with $k_2$)
    \item \bluebox{Unexpand back into modulus}
\end{enumerate}

Just do \yellowbox{these lines} for the first equation, and \greenbox{these lines} for each other equation

% ========== End the document ==========
\end{document}